
\documentclass{article}
\usepackage{colortbl}
\usepackage{makecell}
\usepackage{multirow}
\usepackage{supertabular}

\begin{document}

\newcounter{utterance}

\centering \large Interaction Transcript for game `cladder', experiment `full\_v1.5\_default', episode 1169 with reflexion{-}Qwen3.5{-}27B.
\vspace{24pt}

{ \footnotesize  \setcounter{utterance}{1}
\setlength{\tabcolsep}{0pt}
\begin{supertabular}{c@{$\;$}|p{.15\linewidth}@{}p{.15\linewidth}p{.15\linewidth}p{.15\linewidth}p{.15\linewidth}p{.15\linewidth}}
    \# & \multicolumn{2}{c}{Player} && \multicolumn{2}{c}{Game Master} \\
    \hline

    \theutterance \stepcounter{utterance}  
    & & & \multicolumn{4}{p{0.6\linewidth}}{
        \cellcolor[rgb]{0.9,0.9,0.9}{
            \makecell[{{p{\linewidth}}}]{
                \texttt{\tiny{[P1$\langle$GM]}}
                \texttt{You have attempted this same task before and did not succeed. These are the plans you wrote after those attempts.} \\
\texttt{Trial \#1: New plan: My previous attempt failed because I provided a verbose, multi{-}step derivation and cut off mid{-}sentence, failing to deliver the concise, direct answer required by the prompt's specific formatting constraints ("Start your answer with 'Yes' or 'No', followed by additional reasoning"). I incorrectly assumed a detailed mathematical breakdown was necessary, whereas the task required a direct conclusion based on the calculated conditional probabilities. Next time, I will immediately calculate \$P(R|V)\$ and \$P(R|V\^{}c)\$, compare them, and output the result as a single, complete paragraph starting strictly with "Yes" or "No" followed by a brief summary of the comparison, ensuring the response is not truncated and adheres strictly to the requested structure.} \\
\\ 
\texttt{Trial \#2: New plan: My previous attempt failed because I engaged in visible, unstructured internal monologue and self{-}correction within the final output, leading to a truncated and contradictory response. I explicitly wrote out my confusion ("wait, let me re{-}read") and started a calculation that I never finished. Next time, I will perform all probability calculations (\$P(R|V) = 0.20/0.59 \textbackslash{}approx 0.34\$ and \$P(R|V\^{}c) = 0.24/0.41 \textbackslash{}approx 0.59\$) and logical comparisons internally before generating the response. I will then output a single, coherent paragraph that starts immediately with "Yes" (since \$0.34 < 0.59\$), followed by a concise explanation of the conditional probabilities, ensuring no intermediate thoughts, corrections, or incomplete sentences are included.} \\
\\ 
\\ 
\texttt{You are an expert in causal inference. The following question is not a typical commonsense query, but rather a meticulously designed question created by a professor specializing in causal inference, intended to assess the students' mastery of the course content.} \\
\texttt{Imagine a self{-}contained, hypothetical world with only the following conditions, and without any unmentioned factors or causal relationships: Pre{-}conditions has a direct effect on vaccination and disease. Vaccination has a direct effect on disease. The overall probability of getting the vaccine is 59\%. The probability of vaccine refusal and recovering from the disease is 24\%. The probability of getting the vaccine and recovering from the disease is 20\%. Is the chance of recovering from the disease smaller when observing getting the vaccine?} \\
\texttt{Start your answer with 'Yes' or 'No', followed by additional reasoning or evidence to support your explanation.} \\
            }
        }
    }
     \\ \\

    \theutterance \stepcounter{utterance}  
    & \multicolumn{4}{p{0.6\linewidth}}{
        \cellcolor[rgb]{0.561,0.737,0.561}{
            \makecell[{{p{\linewidth}}}]{
                \texttt{\tiny{[P1$\rangle$GM]}}
                \texttt{No, the chance of recovering from the disease is not smaller when observing getting the vaccine; in fact, it is larger. Based on the provided data, the probability of recovering given vaccination is calculated as \$P(R|V) = P(R \textbackslash{}cap V) / P(V) = 0.20 / 0.59 \textbackslash{}approx 0.339\$. Conversely, the probability of recovering given vaccine refusal (not getting the vaccine) is \$P(R|V\^{}c) = P(R \textbackslash{}cap V\^{}c) / P(V\^{}c) = 0.24 / (1 {-} 0.59) = 0.24 / 0.41 \textbackslash{}approx 0.585\$. Since \$0.339 < 0.585\$, the conditional probability of recovery is actually lower for those who received the vaccine compared to those who refused it in this specific hypothetical scenario, meaning the chance of recovery is *larger* for the unvaccinated group, not smaller for the vaccinated group as the question implies.} \\
            }
        }
    }
    & & \\ \\

    \theutterance \stepcounter{utterance}  
    & & & \multicolumn{2}{p{0.3\linewidth}}{
        \cellcolor[rgb]{0.9,0.9,0.9}{
            \makecell[{{p{\linewidth}}}]{
                \texttt{\tiny{[GM$|$GM]}}
                \texttt{No} \\
            }
        }
    }
    & & \\ \\

    \theutterance \stepcounter{utterance}  
    & & & \multicolumn{2}{p{0.3\linewidth}}{
        \cellcolor[rgb]{0.9,0.9,0.9}{
            \makecell[{{p{\linewidth}}}]{
                \texttt{\tiny{[GM$|$GM]}}
                \texttt{yes} \\
            }
        }
    }
    & & \\ \\

    \theutterance \stepcounter{utterance}  
    & & & \multicolumn{2}{p{0.3\linewidth}}{
        \cellcolor[rgb]{0.9,0.9,0.9}{
            \makecell[{{p{\linewidth}}}]{
                \texttt{\tiny{[GM$|$GM]}}
                \texttt{game\_result = LOSE} \\
            }
        }
    }
    & & \\ \\

    \theutterance \stepcounter{utterance}  
    & & \multicolumn{4}{p{0.6\linewidth}}{
        \cellcolor[rgb]{0.95,0.95,0.95}{
            \makecell[{{p{\linewidth}}}]{
                \texttt{\tiny{[P1$\leftrightarrow$P]}}
                \texttt{Post{-}episode Reflexion (not sent to GM):} \\
\texttt{New plan: My previous attempt failed because I provided a logically contradictory explanation that undermined my correct "No" answer. I correctly calculated that the recovery probability for the vaccinated group (\$\textbackslash{}approx 0.34\$) is lower than for the unvaccinated group (\$\textbackslash{}approx 0.59\$), which means the chance of recovery is indeed smaller when observing vaccination. However, I erroneously stated in the text that the chance was "larger" and that the unvaccinated group had the higher chance, creating a confusing and self{-}defeating argument. Next time, I will strictly align my textual explanation with the mathematical result: I will start with "No," then explicitly state that \$P(R|V) \textbackslash{}approx 0.34\$ is less than \$P(R|V\^{}c) \textbackslash{}approx 0.59\$, and conclude that this confirms the chance of recovery is smaller for the vaccinated group, ensuring the reasoning directly supports the initial "No" without any contradictory phrasing.} \\
            }
        }
    }
    & \\ \\

\end{supertabular}
}

\end{document}
